% Oversells — dizayn namunasi (har bir dizaynning koʻrinishini koʻrsatadi).
% Design showcase: every environment and drawing helper of dizayn.sty.
\documentclass[a4paper]{article}
\usepackage{fontspec}
\setmainfont{Noto Serif}
\usepackage{babel}
\babelprovide[import=uz-Latn, main]{uzbek}
\usepackage{amsmath, amssymb, amsthm}
\usepackage{graphicx}
\usepackage{tikz}
\usepackage{pgfplots}
\usepackage[hidelinks]{hyperref}
\usepackage{dizayn}

\title{Algebra va geometriya}
\author{8-sinf uchun darslik}
\date{2026}

\begin{document}
\maketitle

\section{Kvadrat tenglama}

Ushbu boʻlimda kvadrat tenglamani yechishning umumiy usulini oʻrganamiz.\yon{Diskriminant soʻzi lotincha \emph{discriminare} — «ajratmoq»dan.}
Har bir \belgila{yangi tushuncha} misollar bilan mustahkamlanadi.

\begin{tarif}
  $ax^2 + bx + c = 0$ ($a \neq 0$) koʻrinishidagi tenglama \emph{kvadrat tenglama},
  $D = b^2 - 4ac$ soni esa uning \emph{diskriminanti} deyiladi.
\end{tarif}

\begin{teorema}[Ildizlar formulasi]
  Agar $D \geq 0$ boʻlsa, tenglama haqiqiy ildizlarga ega va ular quyidagicha topiladi:
\end{teorema}

\muhim[Ildizlar formulasi]{x_{1,2} = \frac{-b \pm \sqrt{b^2 - 4ac}}{2a}}

\begin{koordinata}[xmin=-0.5, xmax=5.5, ymin=-1, ymax=6.5]
  \addplot[domain=-0.3:5.3] {x^2 - 5*x + 6};
  \addplot[domain=-0.3:5.3] {0.5*x + 0.5};
  \node[dz nuqta, fill=dzUrgu] at (axis cs:2,0) {};
  \node[dz nuqta, fill=dzUrgu] at (axis cs:3,0) {};
\end{koordinata}

\begin{misol}
  $x^2 - 5x + 6 = 0$ tenglamani yeching.
\end{misol}
\begin{yechim}
  $D = 25 - 24 = 1 > 0$, demak $x_1 = \dfrac{5-1}{2} = 2$, \ $x_2 = \dfrac{5+1}{2} = 3$.
\end{yechim}

\begin{isbot}
  $ax^2+bx+c=0$ ni $4a$ ga koʻpaytirib, toʻliq kvadrat ajratamiz:
  \[ (2ax+b)^2 = b^2-4ac \quad\Longrightarrow\quad 2ax+b=\pm\sqrt{D}. \]
\end{isbot}

\section{Geometriya}

\subsection{Uchburchak va burchaklar}

\begin{chizma}
  \uchburchak{A}{(0,0)}{B}{(5,0)}{C}{(1.6,3)}
  \coordinate (H) at (1.6,0);
  \draw[dz yordamchi] (C) -- (H);
  \togriburchak{C}{H}{B}
  \burchak{B}{A}{C}{\alpha}
  \burchak{C}{B}{A}{\beta}
  \node[below] at (H) {$H$};
\end{chizma}

\begin{teorema}[Pifagor]
  Toʻgʻri burchakli uchburchakda $a^2 + b^2 = c^2$, bunda $c$ — gipotenuza.
\end{teorema}

\begin{eslatma}
  Uchburchak burchaklari yigʻindisi doimo $\alpha + \beta + \gamma = 180^\circ$.
\end{eslatma}

\begin{diqqat}
  $\sqrt{a^2+b^2} \neq a + b$ — bu keng tarqalgan xato!
\end{diqqat}

\subsection{Son oʻqi va kasrlar}

\begin{sonoqi}{-3}{5}
  \oraliq{-1}{3}
  \sonnuqta{-1}{A}
  \sonnuqta{3}{B}
\end{sonoqi}

\noindent Kasrlar: \ \kasrdoira{1}{4}\, $=\dfrac14$, \quad \kasrdoira{3}{8}\, $=\dfrac38$, \quad
\kasrdoira{5}{6}\, $=\dfrac56$. \hfill Toʻplamlar: \venn{A}{B}

\begin{bilasizmi}
  «Algebra» soʻzi Muhammad ibn Muso al-Xorazmiyning «Al-kitob al-muxtasar fi hisob al-jabr
  va-l-muqobala» asari nomidan kelib chiqqan.
\end{bilasizmi}

\begin{mashqlar}
  \item $x^2 - 7x + 12 = 0$ tenglamani yeching. \daraja{1}
  \item $\displaystyle \int_0^1 (3x^2 + 2x)\,dx$ ni hisoblang. \daraja{2}
  \item $\displaystyle \lim_{n\to\infty}\Bigl(1+\frac1n\Bigr)^n = e$ ekanini isbotlang. \daraja{3}
\end{mashqlar}

\begin{javoblar}
  \item $x_1 = 3$, $x_2 = 4$.
  \item $2$.
  \item Bernulli tengsizligidan foydalaning.
\end{javoblar}

\end{document}
